\section{Introduction}

\subsection{Context and Problem Statement}
Legal information is difficult to use when it is scattered across many documents. In Vietnamese labor law, a practical question is often not answered by one paragraph alone. A user may ask about the termination of an employment contract, severance allowance, retirement age, working time, minor workers, or labor disputes. To answer correctly, the system may need to find rules from the Labor Code, a decree, a circular, or even a table or appendix.

This difficulty is increased by the structure of Vietnamese legal documents. Different legal instruments may overlap, refer to each other, provide implementation details, or replace earlier rules~\cite{vietnam2015law80,vietnam2020law63}. A later document may clarify or supersede a previous one, while a lower-level document may explain how a higher-level legal rule should be applied. Therefore, a legal question-answering system should not treat legal texts as independent passages only. It should also preserve the relationships between legal provisions, documents, and legal ranks.

This creates an important challenge for legal question answering. A useful answer should not only be fluent and easy to read. It must also show the legal basis behind the answer. In law, the citation is part of the information itself. Users need to know which document, article, clause, point, or appendix supports the response.

Recent language models can generate natural answers, but they do not automatically guarantee legal grounding~\cite{ji2023survey}. If the relevant legal text is missing, incomplete, outdated, or incorrectly selected, the final answer may still sound convincing. This is a serious limitation for legal information systems. Therefore, the main problem is not only how to generate an answer, but how to retrieve the right legal evidence before generation and check the cited sources after generation.

This thesis addresses that problem through a Vietnamese labor-law AI assistant. The system is designed for citation-grounded question answering over a scoped legal corpus. It retrieves relevant legal provisions, assembles them as supporting evidence, generates an answer from this evidence, and validates whether the final citations are supported by the retrieved context.

\subsection{Objectives and Scope}
The objective of this thesis is to design, implement, and evaluate an AI assistant for Vietnamese labor-law questions, with a case-study focus on employment contract termination and related employment issues.

The first objective is to prepare a structured legal corpus. Instead of treating legal documents as ordinary plain text, the system preserves important legal units such as documents, articles, clauses, points, appendices, and legal rank. This structure helps the system keep track of where each rule comes from.

The second objective is to build a retrieval pipeline for legal evidence. The system combines hybrid retrieval in Qdrant with a Neo4j legal graph. Qdrant is used to search for relevant legal text, while Neo4j is used to represent legal relationships such as hierarchy, cross-references, implementation guidance, and connections between related provisions. This graph-based representation is useful because legal relevance often depends on document relationships, not only textual similarity.

The third objective is to connect retrieval with answer generation. The assistant generates answers from selected legal contexts and then checks whether the cited legal bases actually appear in the retrieved evidence. This step is used to reduce unsupported citations.

The fourth objective is to evaluate the system on a fixed benchmark of 100 questions. The benchmark includes both questions covered by the indexed corpus and questions outside the corpus. This allows the evaluation to measure retrieval performance, answer quality, citation validation, and refusal behavior when there is not enough legal evidence.

The scope of the system is intentionally limited. The indexed corpus includes selected Vietnamese labor-law documents: the Labor Code 2019 and the labor-related subset of the Civil Procedure Code 2015~\cite{vietnam2019laborcode,vietnam2015civilprocedure}, Decree 135/2020/ND-CP and Decree 145/2020/ND-CP~\cite{vietnam2020decree135,vietnam2020decree145}, and Circular 09/2020/TT-BLDTBXH and Circular 10/2020/TT-BLDTBXH~\cite{vietnam2020circular09,vietnam2020circular10}. The system is a legal information assistant for this corpus, not a replacement for professional legal advice.

\subsection{Contributions and Thesis Organization}
The main contributions of this thesis are implementation-oriented. First, it builds a hierarchy-aware processing pipeline for Vietnamese labor-law documents. Second, it implements a retrieval system that combines Qdrant hybrid search with a Neo4j legal graph. Third, it adds citation validation to check whether generated answers remain grounded in retrieved legal evidence. Fourth, it provides a working web-based assistant using a FastAPI backend, a React/Vite frontend, Supabase, Qdrant, Neo4j, and a configured language-model provider. Finally, it evaluates the system using a reproducible 100-question benchmark.

The rest of this thesis is organized as follows. Section~2 reviews the background and related work on retrieval-augmented generation, legal question answering, graph-augmented retrieval, and evaluation for legal QA systems. Section~3 presents the system design and methodology. Section~4 describes the implementation and deployment. Section~5 reports the evaluation results and discusses the main limitations. Section~6 concludes the thesis and suggests future work.
